\section{Matem\'atica cotidiana en C++}

\subsection{Ra\'ices, potencias y logaritmos}

\begin{center}
\begin{tabular}{ll}
Operaci\'on & C++ \\
\hline
$\sqrt x$ real & \texttt{sqrt(x)}, \texttt{sqrtl(x)} \\
$x^y$ real & \texttt{pow(x,y)}, \texttt{powl(x,y)} \\
$e^x,\ln x$ & \texttt{exp(x)}, \texttt{log(x)} \\
$\log_2x,\log_{10}x$ & \texttt{log2(x)}, \texttt{log10(x)} \\
redondeos & \texttt{floor}, \texttt{ceil}, \texttt{round}, \texttt{trunc} \\
hipotenusa estable & \texttt{hypot(x,y)} \\
valor absoluto & \texttt{abs(x)}, \texttt{fabs(x)} \\
\end{tabular}
\end{center}
\warning{
No usar \texttt{pow(a,b)} para calcular potencias enteras exactas, ya que
trabaja con punto flotante y puede introducir errores de redondeo. En su lugar,
usar exponenciación binaria.

Tampoco confiar en convertir directamente \texttt{sqrt(n)} a entero cuando
\texttt{n} está cerca de cuadrados perfectos grandes, debido a posibles errores
de precisión.

Para calcular raíces cuadradas enteras, se recomienda utilizar la función
\texttt{integer\_sqrt}.

Para realizar una división entera redondeando el resultado hacia arriba,
utilizar \texttt{ceil\_division}.

Para potencias enteras usar \texttt{binary\_power}: requiere exponente no
negativo y que la respuesta entre en \texttt{long long}. La sobrecarga real
admite exponentes negativos. Usar \texttt{2LL} o \texttt{2.0L} para elegir
expl\'icitamente la versi\'on entera o real.
}

\begin{lstlisting}[style=cpp]
ll integer_sqrt(ll value) {
    assert(value >= 0);
    ll root = sqrtl((long double)value);
    while ((i128)(root + 1) * (root + 1) <= value) ++root;
    while ((i128)root * root > value) --root;
    return root; // floor(sqrt(value))
}

ll ceil_division(ll a, ll b) {
    assert(b && !(a == numeric_limits<ll>::min() && b == -1));
    ll q = a / b, r = a % b;
    return q + (r != 0 && ((r > 0) == (b > 0)));
}

ll floor_division(ll a, ll b) {
    assert(b && !(a == numeric_limits<ll>::min() && b == -1));
    ll q = a / b, r = a % b;
    return q - (r != 0 && ((r > 0) != (b > 0)));
}

ll binary_power(ll base, ll exponent) {
    assert(exponent >= 0);
    ll result = 1;
    for (; exponent; exponent >>= 1) {
        if (exponent & 1) result *= base;
        if (exponent > 1) base *= base; // Evita el cuadrado final innecesario.
    }
    return result;
}

long double binary_power(long double base, ll exponent) {
    assert(base != 0.0L || exponent >= 0);
    // Convierte |exponent| sin desbordar cuando exponent == LLONG_MIN.
    unsigned long long e = exponent < 0 ? -(unsigned long long)exponent
                                       : (unsigned long long)exponent;
    long double result = 1.0L;
    for (; e; e >>= 1) {
        if (e & 1) result *= base;
        if (e > 1) base *= base;
    }
    return exponent < 0 ? 1.0L / result : result;
}

ll modular_power(ll base, ll exponent, ll modulus) {
    assert(exponent >= 0 && modulus > 0);
    base %= modulus;
    if (base < 0) base += modulus;
    ll result = 1 % modulus;
    for (; exponent; exponent >>= 1) {
        if (exponent & 1) result = (i128)result * base % modulus;
        if (exponent > 1) base = (i128)base * base % modulus;
    }
    return result;
}
\end{lstlisting}

\texttt{integer\_sqrt} ya usa una aproximaci\'on r\'apida con \texttt{sqrtl};
las correcciones enteras conservan exactitud cerca de cuadrados grandes.
Las divisiones requieren divisor no nulo y cociente representable: excluir
\texttt{LLONG\_MIN / -1}. Las potencias usan $O(\log e)$ multiplicaciones
y $O(1)$ memoria; para $e=0$ devuelven la identidad.

\subsubsection{Potencia de 2}

Para potencias de dos, usar desplazamientos.
\warning{El tipo de la izquierda se determina antes de asignar: usar
\texttt{1LL << k} para \texttt{long long}, con $0\le k\le62$ para un resultado
positivo. Para \texttt{unsigned long long}, usar \texttt{1ULL << k} con
$0\le k\le63$. Nunca desplazar por un valor negativo ni por el ancho del tipo
o m\'as.}

\begin{lstlisting} [style=cpp]
    // Potencia 2 elevado a la k; requiere 0 <= k <= 62.
    int k = 40;
    ll x = 1LL << k;

    // Maxima potencia positiva para tipos de 32, 64 y 128 bits.
    int p_int = 1 << 30;
    unsigned int p_uint = 1U << 31;
    ll p_ll = 1LL << 62;
    unsigned long long p_ull = 1ULL << 63;
    i128 p_i128 = (i128)1 << 126;
    u128 p_u128 = (u128)1 << 127;
\end{lstlisting}

\subsection{Punto flotante}

Comparar con error absoluto y relativo:
Ej: 0.1 + 02 == 03?
\begin{lstlisting}[style=cpp]
bool almost_equal(long double a, long double b,
                  long double epsilon = 1e-12L) {
    return fabsl(a - b) <= epsilon * max({1.0L, fabsl(a), fabsl(b)});
}
\end{lstlisting}

\begin{itemize}
  \item \texttt{double}: unas 15 cifras decimales; \texttt{long double}: m\'as en
  GNU/Linux x86, pero no necesariamente en todas las plataformas.
  \item Para $a\cdot b$ antes de m\'odulo, convertir primero a
  \texttt{\_\_int128} si puede exceder \texttt{long long}.
  \item Para $\lceil a/b\rceil$ con positivos sirve
  \texttt{(a+b-1)/b}; con negativos usar la funci\'on anterior.
\end{itemize}

\subsection{Trabajo con \'angulos y trigonometr\'ia}

Usar cuando aparecen tri\'angulos, rotaciones, pendientes, direcciones o
coordenadas polares. C++ recibe y devuelve \'angulos en \textbf{radianes}.

Para un tri\'angulo rect\'angulo:
\[
\sin(\theta)=\frac{CO}{H},\qquad
\cos(\theta)=\frac{CA}{H},\qquad
\tan(\theta)=\frac{CO}{CA}=\frac{\sin(\theta)}{\cos(\theta)}.
\]

El signo depende del cuadrante:
\begin{center}
\begin{tabular}{ccccc}
Cuadrante & rango & $x$ & $y$ & signos positivos \\
\hline
I & $0<\theta<\pi/2$ & $+$ & $+$ & $\sin,\cos,\tan$ \\
II & $\pi/2<\theta<\pi$ & $-$ & $+$ & $\sin$ \\
III & $\pi<\theta<3\pi/2$ & $-$ & $-$ & $\tan$ \\
IV & $3\pi/2<\theta<2\pi$ & $+$ & $-$ & $\cos$ \\
\end{tabular}
\end{center}

\subsubsection{\'Angulos notables}

\begin{center}
\begin{tabular}{ccccc}
Grados & Radianes & $\sin$ & $\cos$ & $\tan$ \\
\hline
$0^\circ$ & $0$ & $0$ & $1$ & $0$ \\
$30^\circ$ & $\pi/6$ & $1/2$ & $\sqrt3/2$ & $1/\sqrt3$ \\
$45^\circ$ & $\pi/4$ & $\sqrt2/2$ & $\sqrt2/2$ & $1$ \\
$60^\circ$ & $\pi/3$ & $\sqrt3/2$ & $1/2$ & $\sqrt3$ \\
$90^\circ$ & $\pi/2$ & $1$ & $0$ & no definida \\
$180^\circ$ & $\pi$ & $0$ & $-1$ & $0$ \\
$270^\circ$ & $3\pi/2$ & $-1$ & $0$ & no definida \\
$360^\circ$ & $2\pi$ & $0$ & $1$ & $0$ \\
\end{tabular}
\end{center}

\subsubsection{Identidades frecuentes}

\begin{align*}
\sin^2\theta+\cos^2\theta &= 1,\\
\sin(-\theta)&=-\sin\theta, & \cos(-\theta)&=\cos\theta,\\
\sin(a+b)&=\sin a\cos b+\cos a\sin b,\\
\cos(a+b)&=\cos a\cos b-\sin a\sin b,\\
\sin(2a)&=2\sin a\cos a,\\
\cos(2a)&=\cos^2a-\sin^2a=1-2\sin^2a.
\end{align*}

La periodicidad es $2\pi$ para seno y coseno, y $\pi$ para tangente.

\begin{center}
\begin{tabular}{lll}
Operaci\'on & \texttt{double} & \texttt{long double} \\
\hline
seno, coseno, tangente & \texttt{sin/cos/tan} & \texttt{sinl/cosl/tanl} \\
inversas & \texttt{asin/acos/atan} & \texttt{asinl/acosl/atanl} \\
\'angulo de $(x,y)$ & \texttt{atan2(y,x)} & \texttt{atan2l(y,x)} \\
hipotenusa & \texttt{hypot(x,y)} & \texttt{hypotl(x,y)} \\
\end{tabular}
\end{center}

\begin{lstlisting}[style=cpp]
const long double PI = acosl(-1.0L);

long double degrees_to_radians(long double degrees) {
    return degrees * PI / 180.0L;
}

long double radians_to_degrees(long double radians) {
    return radians * 180.0L / PI;
}

long double normalize_angle(long double angle) { // [0, 2*pi)
    angle = fmodl(angle, 2.0L * PI);
    if (angle < 0) angle += 2.0L * PI;
    return angle;
}

long double signed_angle(long double angle) { // [-pi, pi)
    angle = normalize_angle(angle);
    if (angle >= PI) angle -= 2.0L * PI;
    return angle;
}

long double safe_acos(long double value) {
    return acosl(clamp(value, -1.0L, 1.0L));
}

long double safe_asin(long double value) {
    return asinl(clamp(value, -1.0L, 1.0L));
}

void trigonometry_example(long double degrees) {
    long double angle = degrees_to_radians(degrees);
    long double sine = sinl(angle);
    long double cosine = cosl(angle);
    long double tangent = tanl(angle);
    long double direction = atan2l(sine, cosine); // primero y, luego x
    (void)tangent; (void)direction;
}
\end{lstlisting}

\subsubsection{Conversi\'on entre grados y radianes}

Una vuelta completa cumple
\[
360^\circ=2\pi\text{ rad},\qquad
180^\circ=\pi\text{ rad}.
\]
Por lo tanto:
\[
rad=grados\frac{\pi}{180},\qquad
grados=rad\frac{180}{\pi}.
\]

En C++17 no existe \texttt{std::numbers::pi}; la forma portable es
\texttt{acosl(-1.0L)}. \texttt{M\_PI} no es est\'andar y puede no estar definido.

\subsubsection{Normalizaci\'on y distancia angular}

El mismo rayo puede representarse con infinitos valores:
$\theta$, $\theta+2\pi$, $\theta-4\pi$, etc. Normalizar evita errores al comparar
direcciones. Para saber cu\'anto girar, no restar \'angulos sin considerar el
salto entre $0$ y $2\pi$.

\begin{lstlisting}[style=cpp]
// Giro antihorario desde 'from' hasta 'to', en [0, 2*pi).
long double counterclockwise_angle(long double from, long double to) {
    return normalize_angle(to - from);
}

// Menor distancia entre dos direcciones, en [0, pi].
long double angle_distance(long double first, long double second) {
    return fabsl(signed_angle(second - first));
}

// Giro firmado mas corto desde 'from' hasta 'to', en [-pi, pi).
long double signed_angle_difference(long double from, long double to) {
    return signed_angle(to - from);
}

// target pertenece al arco antihorario cerrado start -> finish.
bool angle_between_ccw(long double start, long double target,
                       long double finish, long double epsilon = 1e-12L) {
    return counterclockwise_angle(start, target)
        <= counterclockwise_angle(start, finish) + epsilon;
}
\end{lstlisting}

Ejemplo: entre $350^\circ$ y $10^\circ$ hay solamente $20^\circ$, no
$340^\circ$.

\subsubsection{\texttt{atan2}: direcci\'on y cuadrantes}

Para un vector $(x,y)$:
\[
r=\sqrt{x^2+y^2},\qquad \theta=\operatorname{atan2}(y,x).
\]

Siempre escribir \texttt{atan2(y,x)}: el argumento vertical va primero.
\texttt{atan(y/x)} pierde el cuadrante, divide por cero cuando $x=0$ y no puede
distinguir $(1,1)$ de $(-1,-1)$.

\begin{center}
\begin{tabular}{ccc}
Vector & \texttt{atan2(y,x)} & grados \\
\hline
$(1,0)$ & $0$ & $0^\circ$ \\
$(0,1)$ & $\pi/2$ & $90^\circ$ \\
$(-1,0)$ & $\pi$ & $180^\circ$ \\
$(0,-1)$ & $-\pi/2$ & $-90^\circ$ \\
\end{tabular}
\end{center}

\warning{La direcci\'on del vector $(0,0)$ no est\'a definida. Algunas
implementaciones devuelven cero para \texttt{atan2(0,0)}, pero no debe usarse
como un \'angulo geom\'etrico v\'alido.}

\subsubsection{Coordenadas polares y rotaciones}

Conversi\'on polar a cartesiana:
\[
x=r\cos\theta,\qquad y=r\sin\theta.
\]
Rotar $(x,y)$ un \'angulo $\theta$ en sentido antihorario:
\[
x'=x\cos\theta-y\sin\theta,\qquad
y'=x\sin\theta+y\cos\theta.
\]

\begin{lstlisting}[style=cpp]
pair<long double,long double> polar_to_cartesian(
        long double radius, long double angle) {
    return {radius * cosl(angle), radius * sinl(angle)};
}

pair<long double,long double> rotate_coordinates(
        long double x, long double y, long double angle) {
    long double cosine = cosl(angle), sine = sinl(angle);
    return {x * cosine - y * sine,
            x * sine + y * cosine};
}

pair<long double,long double> cartesian_to_polar(
        long double x, long double y) {
    return {hypotl(x, y), atan2l(y, x)}; // {radio, angulo}
}
\end{lstlisting}

Para rotar alrededor de un centro $C$, trasladar $P-C$, rotar y volver a sumar
$C$. Varias rotaciones suman sus \'angulos; una rotaci\'on no cambia distancias.

\subsubsection{\'Angulo entre vectores}

Para vectores $u$ y $v$ no nulos:
\[
\cos\theta=\frac{u\cdot v}{|u||v|},\qquad
\theta=\operatorname{atan2}(|u\times v|,u\cdot v).
\]

La versi\'on con \texttt{atan2} suele ser m\'as estable y conserva mejor
\'angulos cercanos a $0$ o $\pi$.

\begin{lstlisting}[style=cpp]
long double angle_between_vectors(long double ux, long double uy,
                                  long double vx, long double vy) {
    long double dot = ux * vx + uy * vy;
    long double cross_value = ux * vy - uy * vx;
    assert(hypotl(ux, uy) > 0 && hypotl(vx, vy) > 0);
    return atan2l(fabsl(cross_value), dot); // [0, pi]
}

// Angulo orientado u -> v: positivo antihorario, en [-pi, pi].
long double oriented_angle(long double ux, long double uy,
                           long double vx, long double vy) {
    return atan2l(ux * vy - uy * vx, ux * vx + uy * vy);
}
\end{lstlisting}

Si s\'olo importa comparar \'angulos o saber si son agudos/obtusos, evitar la
trigonometr\'ia: $u\cdot v>0$ indica \'angulo agudo, $=0$ recto y $<0$ obtuso.

\subsubsection{Tri\'angulos cualesquiera}

Si los lados $a,b,c$ est\'an opuestos a los \'angulos $A,B,C$:

\textbf{Ley de cosenos:}
\[
c^2=a^2+b^2-2ab\cos C,\qquad
\cos C=\frac{a^2+b^2-c^2}{2ab}.
\]

\textbf{Ley de senos:}
\[
\frac{a}{\sin A}=\frac{b}{\sin B}=\frac{c}{\sin C}=2R,
\]
donde $R$ es el radio de la circunferencia circunscrita.

\textbf{\'Area:}
\[
s=\frac{a+b+c}{2},\qquad
K=\sqrt{s(s-a)(s-b)(s-c)}.
\]

\begin{lstlisting}[style=cpp]
long double angle_opposite_side(long double a, long double b,
                                long double opposite) {
    assert(a > 0 && b > 0);
    long double cosine = (a * a + b * b - opposite * opposite)
                       / (2.0L * a * b);
    return safe_acos(cosine);
}

long double opposite_side(long double a, long double b,
                          long double included_angle) {
    long double squared = a * a + b * b
                        - 2.0L * a * b * cosl(included_angle);
    return sqrtl(max(0.0L, squared));
}

bool valid_triangle(long double a, long double b, long double c,
                    long double epsilon = 1e-12L) {
    return a > 0 && b > 0 && c > 0
        && a + b > c + epsilon
        && a + c > b + epsilon
        && b + c > a + epsilon;
}

long double triangle_area_heron(long double a, long double b,
                                long double c) {
    assert(valid_triangle(a, b, c));
    long double semiperimeter = (a + b + c) / 2.0L;
    long double product = semiperimeter * (semiperimeter - a)
                        * (semiperimeter - b) * (semiperimeter - c);
    return sqrtl(max(0.0L, product));
}
\end{lstlisting}

\warning{La ley de senos puede tener dos soluciones cuando se conocen dos lados
y un \'angulo no comprendido. La ley de cosenos determina directamente el lado
opuesto al \'angulo comprendido. Her\'on puede perder precisi\'on en tri\'angulos
casi degenerados; con coordenadas, preferir producto cruzado.}

\subsubsection{Circunferencias, arcos y sectores}

Con radio $r$ y \'angulo central $\theta$ en radianes:
\[
\text{arco}=r\theta,\qquad
\text{sector}=\frac12r^2\theta,\qquad
\text{cuerda}=2r\sin(\theta/2).
\]

\begin{lstlisting}[style=cpp]
long double arc_length(long double radius, long double angle) {
    return radius * fabsl(angle);
}

long double sector_area(long double radius, long double angle) {
    return radius * radius * fabsl(angle) / 2.0L;
}

long double chord_length(long double radius, long double angle) {
    return 2.0L * radius * sinl(fabsl(angle) / 2.0L);
}
\end{lstlisting}

\subsubsection{Pol\'igonos regulares}

Para un pol\'igono regular de $n\ge3$ lados:
\[
\sum\text{ interiores}=(n-2)\pi,\quad
\text{interior}=\frac{(n-2)\pi}{n},\quad
\text{exterior}=\frac{2\pi}{n}.
\]
Si $R$ es el circunradio y $s$ el lado:
\[
s=2R\sin(\pi/n),\qquad
\text{\'area}=\frac{nR^2\sin(2\pi/n)}2.
\]

\subsubsection{Convenciones: matem\'atica, br\'ujula y pantalla}

Antes de usar un \'angulo, identificar la convenci\'on del problema:
\begin{itemize}
  \item matem\'atica: $0$ apunta a $+x$ y crece antihorario;
  \item rumbo de br\'ujula: $0$ apunta al norte y crece horario;
  \item matrices/pantallas: la fila o coordenada $y$ suele crecer hacia abajo,
  invirtiendo visualmente la orientaci\'on.
\end{itemize}

\begin{lstlisting}[style=cpp]
long double bearing_to_math_angle(long double bearing) {
    return normalize_angle(PI / 2.0L - bearing);
}

long double math_angle_to_bearing(long double angle) {
    return normalize_angle(PI / 2.0L - angle);
}

// Direccion en una grilla donde row crece hacia abajo.
long double screen_direction(long double delta_column,
                             long double delta_row) {
    return atan2l(-delta_row, delta_column);
}

// Una recta no orientada es igual tras girar pi: resultado [0, pi).
long double normalize_line_angle(long double angle) {
    angle = fmodl(angle, PI);
    if (angle < 0) angle += PI;
    return angle;
}
\end{lstlisting}

\subsubsection{Orden polar exacto sin funciones trigonom\'etricas}

Para ordenar muchos vectores por \'angulo, calcular \texttt{atan2} para cada uno
es m\'as lento y agrega error flotante. Se puede separar el plano en dos mitades
y comparar con producto cruzado. El siguiente orden comienza en $+x$ y avanza
antihorario; los colineales se ordenan por longitud.

\begin{lstlisting}[style=cpp]
int polar_half(ll x, ll y) {
    return (y < 0 || (y == 0 && x < 0)) ? 1 : 0;
}

bool polar_less(pair<ll,ll> first, pair<ll,ll> second) {
    auto [ax, ay] = first;
    auto [bx, by] = second;
    int half_a = polar_half(ax, ay), half_b = polar_half(bx, by);
    if (half_a != half_b) return half_a < half_b;
    i128 cross_value = (i128)ax * by - (i128)ay * bx;
    if (cross_value != 0) return cross_value > 0;
    i128 length_a = (i128)ax * ax + (i128)ay * ay;
    i128 length_b = (i128)bx * bx + (i128)by * by;
    return length_a < length_b;
}
\end{lstlisting}

\warning{Excluir el vector $(0,0)$ o decidir expl\'icitamente d\'onde ubicarlo:
no posee direcci\'on. El comparador debe ser estricto; no introducir epsilon en
un comparador de \texttt{sort}, \texttt{set} o \texttt{map}.}

\subsubsection{Proyecci\'on, reflexi\'on y bisectriz}

La proyecci\'on del vector $v$ sobre una direcci\'on no nula $d$ es
\[
\operatorname{proj}_d(v)=d\frac{v\cdot d}{d\cdot d}.
\]
Si $n$ es una normal de una recta/espejo, la reflexi\'on es
\[
\operatorname{reflect}(v)=v-2\operatorname{proj}_n(v).
\]
La bisectriz interna de dos vectores no opuestos apunta en la suma de sus
versiones unitarias.

\begin{lstlisting}[style=cpp]
pair<long double,long double> project_vector(
        long double vx, long double vy,
        long double dx, long double dy) {
    long double norm_squared = dx * dx + dy * dy;
    assert(norm_squared > 0);
    long double scale = (vx * dx + vy * dy) / norm_squared;
    return {dx * scale, dy * scale};
}

pair<long double,long double> reflect_vector(
        long double vx, long double vy,
        long double normal_x, long double normal_y) {
    auto projection = project_vector(vx, vy, normal_x, normal_y);
    return {vx - 2.0L * projection.first,
            vy - 2.0L * projection.second};
}

optional<pair<long double,long double>> internal_angle_bisector(
        long double ax, long double ay,
        long double bx, long double by) {
    long double norm_a = hypotl(ax, ay), norm_b = hypotl(bx, by);
    if (norm_a == 0 || norm_b == 0) return nullopt;
    long double x = ax / norm_a + bx / norm_b;
    long double y = ay / norm_a + by / norm_b;
    long double norm = hypotl(x, y);
    if (norm < 1e-18L) return nullopt; // direcciones opuestas
    return pair<long double,long double>{x / norm, y / norm};
}
\end{lstlisting}

Para reflejar un punto respecto de una recta, proyectar el vector desde un punto
de la recta hasta el punto original y usar la misma identidad. Para la bisectriz
externa, restar los vectores unitarios en lugar de sumarlos.

\subsubsection{Funciones relacionadas de \texttt{<cmath>}}

\begin{itemize}
  \item \texttt{sinh/cosh/tanh}: funciones hiperb\'olicas; no describen \'angulos
  circulares y aparecen ocasionalmente en f\'ormulas anal\'iticas.
  \item \texttt{fmod(a,b)} conserva el signo de $a$; para normalizar \'angulos
  negativos hay que sumar el per\'iodo, como hace \texttt{normalize\_angle}.
  \item \texttt{remainder(a,b)} devuelve el residuo respecto del m\'ultiplo m\'as
  cercano y puede ser \'util para obtener un \'angulo firmado, aunque su tratamiento
  de empates suele ser menos intuitivo en competencia.
  \item \texttt{copysign(magnitude,sign)} copia el signo sin bifurcar;
  \texttt{isfinite/isnan/isinf} permiten detectar resultados inv\'alidos.
\end{itemize}

\textbf{Rangos de retorno:}
\begin{itemize}
  \item \texttt{asin}: $[-\pi/2,\pi/2]$; \texttt{acos}: $[0,\pi]$.
  \item \texttt{atan}: $(-\pi/2,\pi/2)$; \texttt{atan2(y,x)}: $[-\pi,\pi]$.
  \item \texttt{atan2} distingue los cuatro cuadrantes y funciona cuando
  $x=0$; para direcciones casi siempre es mejor que \texttt{atan(y/x)}.
\end{itemize}

\warning{Por errores de redondeo, un coseno calculado puede quedar apenas fuera
de $[-1,1]$ y hacer que \texttt{acos} devuelva NaN; aplicar \texttt{clamp}.
La tangente no est\'a definida cuando $\cos(\theta)=0$.}

\textbf{Checklist de competencia:}
\begin{itemize}
  \item confirmar si la entrada est\'a en grados o radianes;
  \item usar \texttt{atan2(y,x)}, no \texttt{atan(y/x)};
  \item normalizar antes de comparar direcciones que cruzan $0/2\pi$;
  \item aplicar \texttt{clamp} antes de \texttt{asin/acos};
  \item evitar trigonometr\'ia si dot/cross permiten decidir exactamente;
  \item usar \texttt{long double} y un epsilon acorde a las coordenadas;
  \item no comparar resultados flotantes con \texttt{==}.
\end{itemize}

Cada funci\'on trigonom\'etrica y cada helper de esta secci\'on cuesta $O(1)$
tiempo y $O(1)$ memoria.
